\documentclass[11pt,a4paper]{article}
\usepackage{iftex}
\usepackage{amsmath,amsthm}
\ifPDFTeX
  \usepackage[T1]{fontenc}
  \usepackage[utf8]{inputenc}
  \usepackage{textcomp}
  \usepackage{amssymb}
\else
  % XeTeX/LuaTeX: amssymb BEFORE unicode-math (unicode-math then overrides it cleanly;
  % the reverse order clashes on \eth). Keeps \square, \mathbb, \mathfrak available.
  \usepackage{amssymb}
  \usepackage{unicode-math}
  \defaultfontfeatures{Scale=MatchLowercase}
\fi
\usepackage{booktabs}
\usepackage{longtable}
\usepackage{array}
\usepackage{graphicx}
\usepackage[margin=1in]{geometry}
\usepackage{xcolor}
\usepackage{hyperref}
\hypersetup{colorlinks=true,linkcolor=blue!60!black,citecolor=green!50!black,urlcolor=blue!60!black}
\usepackage{fancyhdr}
\pagestyle{fancy}
\fancyhf{}
\fancyhead[L]{\small PACSeries Paper 12}
\fancyhead[R]{\small Dawn Field Institute}
\fancyfoot[C]{\thepage}
\renewcommand{\headrulewidth}{0.4pt}
\setlength{\parindent}{0pt}
\setlength{\parskip}{6pt plus 2pt minus 1pt}
\providecommand{\tightlist}{\setlength{\itemsep}{0pt}\setlength{\parskip}{0pt}}

\title{First Observational Contact: Testing the Cascade Clock Against Quasar Absorption Lines}
\author{Peter Groom \\ Dawn Field Institute}
\date{June 2026}

\begin{document}
\maketitle

\subsection*{A test of the cascade clock's universality against quasar absorption lines}

\textbf{Peter Groom, Dawn Field Institute} \\
\textbf{PACSeries Paper 12} \\
\textbf{Date}: June 2026 \\
\textbf{Version}: 1.0 (Draft)

\begin{center}\rule{0.5\linewidth}{0.4pt}\end{center}

\section*{Abstract}

We make first observational contact between the DFT cascade clock and quasar absorption spectroscopy. The clock $N(z) = 1.360 + (1/\ln\varphi) \cdot \ln(t_{\text{lookback}})$, previously validated on three cosmological observables (S8, H0, JWST), generates one sharp, DFT-specific prediction for absorption lines: non-monotonic width oscillation at integer cascade levels, with named transition redshifts $z = 0.101, 0.171, 0.302, 0.579, 1.416$. This prediction was pre-registered on GitHub (commit \texttt{193d1c8e}) before any observational data was examined.

The oscillation prediction is falsified. Z-detrending kills every oscillatory signal. The surviving smooth correlation — $b \propto \ln t_{\text{lookback}}$ at $R^2 = 0.85$ over 443,000 CIV systems — is consistent with the predicted logarithmic form, with the $\varphi$-determined slope costing zero $R^2$ against the free fit, and it outperforms the standard halo-virial scaling, which collapses to $\alpha \approx 0$ when fit ($\Delta\text{BIC} = 44$). But it remains degenerate with generic cosmic-time evolution ($z^2$ achieves $R^2 = 0.86$ with one more parameter). Absorption spectroscopy in the well-sampled $z = 1.5$–$4.5$ range does not have the leverage to discriminate the cascade clock from its smooth mimics.

The paper contributes three things. First, a pre-registered, falsifiable prediction. Its falsification localizes where the clock loses discriminating power and establishes a reusable z-detrending methodology. Second, a two-channel concept: DFT predicts photons carry a conserved channel (PAC: line ratios, invariant to machine precision via integer ADE topology) and a historical channel (SEC: line profiles, epoch-dependent). This partition survives regardless of the cascade clock's parameterization and offers a new separation of structural from thermodynamic information in multi-line spectroscopy. Third, a forward program: the cascade clock and its polynomial mimics diverge outside the fitted range, making $z > 5$ (JWST) and $z < 0.5$ (CaII) the discriminating tests.

\textbf{Keywords}: cascade clock, quasar absorption lines, falsification, pre-registration, PAC conservation, Dawn Field Theory

\textbf{Epistemic status.} This paper is part of the PACSeries, an evolving research program rather than a finished theory (see the series note \emph{How to Read the PACSeries}). Its quantitative claims are tiered by derivation type — \textbf{A} structural, \textbf{B} identified, \textbf{C} pattern-matched — and labeled in the text. Failures are reported rather than omitted, and framework-level corrections are logged in the Epistemic Corrections Registry in the repository. Every headline result is reproducible from the \texttt{Code/} and \texttt{Data/} in this package (provenance in \texttt{trace.yaml}).

\begin{center}\rule{0.5\linewidth}{0.4pt}\end{center}

\section*{1. The question and the test}

\subsection*{1.1 What the cascade clock predicts}

The DFT cascade clock (Paper 9) derives from the PAC conservation axiom. The temporal function $N(t) = a + (1/\ln\varphi) \cdot \ln(t_{\text{lookback}})$ was calibrated on three independent cosmological data points:

\begin{itemize}\tightlist
\item $S_8(z = 0.35) = 0.769$ vs $0.768$ observed ($3.22\sigma \to 0.07\sigma$)
\item $H_0$ ratio $= \varphi^{1/6}$ vs $1.0838$ observed ($0.076\%$)
\item JWST galaxy counts at $z = 8$ and $z = 12$
\end{itemize}

The question: does this same clock, with no additional parameters, predict the evolution of individual ionic transitions in galaxy-halo gas clouds?

\subsection*{1.2 The discriminating prediction}

We derived a DFT-specific prediction: at integer cascade levels $N$, the PAC ledger is transitioning. Line widths should peak. At half-integer $N$, the ledger is settled. Line widths should narrow. This produces a non-monotonic oscillation with specific transition redshifts determined by the clock parameters. This prediction was registered on GitHub (commit \texttt{193d1c8e}, pushed June 6, 2026) before any SDSS data was examined.

\subsection*{1.3 Degeneracy note}

$N(z)$ is a monotone function of redshift. Any observable that evolves monotonically with cosmic time will correlate with it. Smooth correlations with $N$ are therefore necessary but not sufficient evidence for the cascade clock. Discriminating power requires either non-monotonic structure (the oscillation) or extrapolation outside the fitted range where the log form and polynomial mimics diverge.

\begin{center}\rule{0.5\linewidth}{0.4pt}\end{center}

\section*{2. Data}

\textbf{SDSS DR16 MgII}: 89,291 systems ($z = 0.35$–$2.28$) from MPA-Garching. \\
\textbf{SDSS DR16 FeII-confirmed}: 52,872 systems with 4 transitions. \\
\textbf{SDSS DR12 CIV}: 443,000 systems ($z = 1.4$–$5.0$) from Monadi et al. (2023). \\
\textbf{XQR-30}: 5,764 components, 42 sightlines ($z = 2.0$–$6.5$), 8 ionic species (Davies et al. 2023).

\begin{center}\rule{0.5\linewidth}{0.4pt}\end{center}

\section*{3. The oscillation prediction: falsified}

\subsection*{3.1 Initial signals and their collapse}

Our first experiments found cascade-correlated signals: EW spread ($p = 0.007$), doublet coupling ($p = 0.006$), population separation ($p = 10^{-42}$). Z-detrending — removing a quadratic fit in $z$ — killed every one:

\begin{longtable}[]{@{}lll@{}}
\toprule
Signal & Raw $p$ & Detrended $p$ \\
\midrule\endhead
MgII EW spread & 0.007 & 0.87 \\
Doublet FWHM discrepancy & 0.006 & 0.47 \\
CIV doublet ratio & $\approx 0$ & 0.68 \\
\bottomrule
\end{longtable}

The oscillatory component is absent from the data.

\subsection*{3.2 What the falsification means}

The clock's non-monotonic feature does not imprint on absorption line statistics at this precision. This does not falsify the cascade clock itself — it remains validated on S8/H0/JWST — but it falsifies the specific prediction that SEC restructuring at cascade boundaries produces detectable oscillations.

\begin{center}\rule{0.5\linewidth}{0.4pt}\end{center}

\section*{4. The smooth survivor}

\subsection*{4.1 CIV velocity tracks lookback time}

The CIV Doppler $b$-parameter (binned medians, 97 bins over $z = 1.5$–$4.5$, 443,642 systems) correlates with $\ln(t_{\text{lookback}})$ at $R^2 = 0.853$. The data is consistent with the predicted logarithmic form. A second observable from the same catalog — the CIV doublet equivalent-width ratio — shows the same logarithmic preference at lower amplitude ($R^2 = 0.51$ vs $0.45$ for linear $z$; exp\_12). The two observables are not conflated in what follows: all model-comparison numbers refer to the $b$-parameter (exp\_18).

Fixing the slope to the $\varphi$-determined value $1/\ln\varphi = 2.0781$ — rather than fitting it — costs zero $R^2$: the constrained-to-free slope ratio is $1.000000$ at machine precision (exp\_12 T2, exp\_18). The data is exactly as consistent with $\varphi$-rate evolution as with the best free logarithmic fit. This is the strongest pro-clock statement the interpolation range supports, with the caveat of §4.2.

\subsection*{4.2 Degeneracy limits}

The zero-cost result does not test $\varphi$ specifically: $N(z)$ is an affine function of $\ln(t)$, so the free per-observable amplitude absorbs the $\varphi$-determined slope. Against the correct null — smooth functions of cosmic time — the cascade clock beats the standard astrophysical model but does not uniquely win (exp\_18; BIC = $n\ln(\text{SS}_{\text{res}}/n) + k\ln n$, lower is better):

\begin{longtable}[]{@{}llll@{}}
\toprule
Model & Parameters & $R^2$ & BIC \\
\midrule\endhead
$z$ (linear) & 2 & 0.720 & 387.8 \\
Halo virial $A + B(1+z)^{\alpha}$ & 3 & 0.780 & 368.8 \\
$\ln(t)$ / cascade clock & 2 & 0.853 & 325.2 \\
$z^2$ (quadratic) & 3 & 0.864 & 322.2 \\
$z^3$ (cubic) & 4 & 0.877 & 317.1 \\
\bottomrule
\end{longtable}

Two readings. Against standard astrophysics the clock wins outright: the halo virial scaling — gas velocity tracking $(1+z)^{1/2}$ at fixed halo mass — collapses to $\alpha = 0.0003$ when fit, and trails the clock by $\Delta\text{BIC} = 44$ despite an extra parameter. Against generic smooth mimics the clock ties or loses: $z^2$ edges it by $\Delta\text{BIC} = 3$ with one more parameter, and the cubic by $8$ with two. The data confirms that CIV velocity evolves smoothly with cosmic time in a way standard halo physics does not capture — consistent with the cascade clock but not uniquely selected by it.

\subsection*{4.3 Supporting trends}

Three additional trends are consistent with smooth cosmic-time evolution:

\textbf{Velocity skewness} transitions from symmetric (high $N$, early) to right-skewed (low $N$, late) with $\rho = -0.929$, $p = 0.003$ over 7 bins. This is consistent with the cascade's prediction of higher information-processing rate at earlier epochs, but also with standard expectations: earlier gas is more turbulent (symmetric) and later gas is more structured (skewed). The correlation is over 7 binned medians.

\textbf{Fe/Mg ratio} decreases with $N$ ($\rho = -0.949$, $R^2 = 0.89$, 19 bins), matching the nucleosynthesis timescale. This is textbook chemical evolution — less Type Ia iron at earlier times — reproduced by any monotonic clock.

\textbf{Ionization redistribution}: 8 ions in XQR-30 show low-ionization species weakening and high-ionization strengthening with $N$ (FeII $p = 0.000$, SiII $p = 0.001$, SiIV $p = 0.011$, CIV $p = 0.000$). This is consistent with UV-background hardening, the standard explanation for cosmic ionization evolution.

All three are \textbf{[C]}: pattern-consistent with the cascade clock, also explained by standard astrophysics.

\subsection*{4.4 Signals that survive detrending}

Three findings are immune to the z-detrending that killed the oscillation — not because they survived it, but because they are constructed so that a smooth z-trend cannot produce them:

\textbf{Sightline-straddling pair coherence} (exp\_08, panel C): inter-absorber correlations between systems on the same sightline that straddle a cascade transition redshift, compared against same-separation pairs that do not. 6,437 straddling vs 9,000 control pairs; Mann-Whitney $p = 1.8 \times 10^{-9}$, KS $p = 1.7 \times 10^{-10}$. A smooth trend affects both pair classes identically.

\textbf{Narrow-window doublet coherence} (exp\_08, panel D): doublet-ratio distribution shape inside narrow windows at transition redshifts versus troughs. 14,020 transition vs 10,235 trough systems; KS $p = 1.3 \times 10^{-5}$ (doublet ratio), $p = 7.5 \times 10^{-5}$ (width discrepancy). Within a narrow window the z-trend is locally constant.

\textbf{Entropy gradients along sightlines} (exp\_10 panel C, exp\_11 follow-up): 557 of 1,620 multi-absorber sightlines (34\%) show monotonic entropy ordering, roughly twice the random expectation.

These are classified \textbf{[B]} at best and reported with deliberate caution: each appeared in a single experimental pass, none has been independently replicated, and the regularity experiments that followed (exp\_15–17) tested other channels and did not retest them. They are the strongest candidates for cascade structure beyond smooth evolution, and the first targets for an adversarial replication pass before any is promoted to a claim.

\begin{center}\rule{0.5\linewidth}{0.4pt}\end{center}

\section*{5. The two-channel concept}

\subsection*{5.1 The partition}

DFT predicts photons carry two channels:

\textbf{PAC channel (conserved)}: spectral line ratios, determined by discrete ADE graph topology. The adjacency matrix is integer-valued and cannot evolve continuously. The conserved class is invariant exactly, to machine precision, because the underlying object is a graph with integer entries. This topological exactness is DFT-specific: the Standard Model also predicts invariant ratios, but DFT predicts \emph{why} (integer topology) and \emph{how precisely} (exactly). The derivation of $\alpha_{\text{EM}}$ to 5.7 ppm from Fibonacci structure (Paper 4, \textbf{[A]}) is a hard commitment: $\alpha$ cannot drift because $\varphi$ is a fixed point and Fibonacci numbers are integers. Confirming $\alpha$-invariance does not discriminate DFT from the SM (both predict it), but violation would falsify DFT, an asymmetric test.

\textbf{SEC channel (historical)}: spectral line widths, shapes, and profiles, determined by the entropy state at the absorption site. Different epochs have different SEC states; the profiles evolve. The cascade clock parameterizes this evolution, though not uniquely (§4.2).

\subsection*{5.2 What the partition enables}

The novel claim is the partition itself: DFT tells you \emph{a priori} which observables encode structure (conserved, topological) and which record history (evolving, thermodynamic). Current absorption-line analysis measures metallicity and kinematics without distinguishing which information is structural and which is historical. The PAC/SEC partition offers this separation. If correct, every multi-line absorption system carries more extractable information than standard analysis recovers.

\begin{center}\rule{0.5\linewidth}{0.4pt}\end{center}

\section*{6. What would discriminate}

\subsection*{6.1 Extrapolation, not interpolation}

Within $z = 1.5$–$4.5$, every smooth function fits. Outside that range, $\ln(t)$ and polynomial mimics diverge:

\begin{itemize}\tightlist
\item At $z > 5$: the log form flattens while polynomials extrapolate. JWST NIRSpec data at $z = 5$–$7$ is the kill test.
\item At $z < 0.5$: the log form steepens sharply — $dN/dz = 12.4$ at $z = 0.15$ vs $2.9$ at $z = 0.5$, a $4.3\times$ gradient change that no CIV-trained polynomial reproduces.
\end{itemize}

The low-z test was performed (exp\_18). A discrimination test against the 435 SDSS CaII absorbers of Sardane et al. (2014) ($z = 0.03$–$1.34$, $\lambda 3934$ rest EW) was pre-registered — locked clock, binning rule, and $\Delta\text{BIC} = \pm 6$ decision thresholds — at commit \texttt{fbad01d1}, before the catalog was downloaded. The registered verdict is \textbf{inconclusive}: $\Delta\text{BIC}(\text{clock} - \text{best mimic}) = +0.3$ in both registered tests, far inside the thresholds. The informative part is why every model fails: the best fit anywhere is a 4-parameter cubic at $R^2 = 0.16$, with all 2-parameter forms at $R^2 \approx 0.03$–$0.05$. CaII median EW carries essentially no redshift trend in this sample, consistent with the registered threat that CaII selects dusty sightlines whose statistics are local, not cosmological. The low-z edge cannot discriminate the clock with an EW observable at this sample size; a kinematic observable (Doppler $b$ at $z < 0.5$) is what the test needs.

\textbf{Prediction}: the cascade clock's discriminating power lives at the edges of the current data, not in the middle. With the CaII EW channel exhausted, the burden rests on $z > 5$ (JWST) and low-z kinematics.

\subsection*{6.2 Cross-observable shared normalization (performed)}

We tested whether CIV $b$, MgII FWHM, and Fe/Mg ratio share a common shape against $N(z)$. They do not: CIV increases (slope $+92$), MgII decreases ($-2.6$), Fe/Mg decreases ($-0.04$). The anti-correlation ($r = -0.63$ between CIV and Fe/Mg) is the ionization redistribution plane in another form. The shapes are complementary, not shared, and do not break the degeneracy.

\subsection*{6.3 A predicted crossover energy: derived, registered, first measurement attempted}

The ionization redistribution crossover lies between AlIII (18.8 eV) and SiIV (33.5 eV). UV-background hardening does not predict the crossover location; a derived value would be DFT-specific.

The derivation now exists. The Milestone R energy-scale machinery (exp\_24: EM-scale energies are $\alpha(d)^2 m_{\text{mediator}}$, validated by the Rydberg at 11.4 ppm) gives the ledger-severance cost as the full Coulomb energy $\alpha^2 m_e c^2$ = one Hartree = 27.2 eV, at the center of the observed bracket. The prediction, including the factor-of-2 argument (severance costs the full interaction energy, not the virial-halved binding energy), was registered at commit \texttt{d9c77a81} before measurement.

The first measurement (exp\_19, pre-registered metric: per-ion fractional EW slope vs $N(z)$ across seven ions in four catalogs) returned \textbf{inconclusive}: the coupling metric is not portable across surveys with different selection and EW regimes, and the cross-survey curve is non-monotone by construction. Within homogeneous data the law's structure appears: CaII coupling is consistent with zero ($\beta = +0.03$, CI95 $[-0.06, +0.10]$ — the settled-phase anchor, now with an error bar), the SDSS-only ordering rises monotonically (CaII $\to$ FeII $\to$ MgII $\to$ CIV), and XQR-30 alone reproduces the sign flip between 8.2 and 33.5 eV. Localizing the zero crossing — the actual test of 27.2 eV — requires intermediate-IP statistics (CII at 11.3 eV, AlIII at 18.8 eV) that current public catalogs do not provide at the registered quality floor. The Hartree prediction stands, untested, for that dataset.

\begin{center}\rule{0.5\linewidth}{0.4pt}\end{center}

\section*{7. Failure inventory}

\begin{longtable}[]{@{}llll@{}}
\toprule
\# & Prediction & Result & Lesson \\
\midrule\endhead
1 & Width oscillation at integer $N$ & Falsified & Cascade produces smooth evolution, not oscillation \\
2–4 & Binned cascade correlations & z-trend confounds & z-detrending is mandatory for cascade claims \\
5 & N-space periodicity & Not detected & No cascade-frequency power \\
6 & Sharp excess at transition $z$ & Not above controls & No localized features \\
7 & Single-tree rotation curves & Failed & Network model needed \\
8 & Cosmic velocity $\approx$ lab turbulence & Anti-correlated & Different coupling regime \\
\bottomrule
\end{longtable}

Two of these deserve a sentence beyond the table. Entry 7: modeling a galaxy halo as a single PAC tree fails to reproduce flat rotation curves — the cascade ledger of one tree does not carry enough structure, and any future attempt needs coupled trees (a network) rather than a deeper single cascade. Entry 8: laboratory turbulence velocity statistics \emph{anti}-correlate with the cosmic velocity evolution rather than matching it, indicating the absorber gas sits in a different coupling regime than tabletop flows — the M2/M4 turbulence results do not transfer directly to cosmological gas.

\begin{center}\rule{0.5\linewidth}{0.4pt}\end{center}

\section*{8. Classification}

\begin{longtable}[]{@{}lll@{}}
\toprule
Finding & Class & Justification \\
\midrule\endhead
$\alpha$ invariance & \textbf{[A]} & Derived commitment; kills DFT if violated; confirmation doesn't discriminate from SM \\
Integer-$N$ oscillation & \textbf{[A] falsified} & Derived prediction; killed by z-detrending \\
Topological exactness of PAC channel & \textbf{[B]} & Novel mechanism (integer ADE adjacency) \\
Two-channel partition & \textbf{[B]} & Derived + identified with line ratios/profiles \\
$b \propto \ln(t)$ & \textbf{[C]} & Consistent; degenerate with cosmic time \\
$\varphi$-slope zero cost & \textbf{[C]} & Exact consistency, but amplitude freedom absorbs the test \\
Halo virial collapse ($\alpha \approx 0$) & \textbf{[B]} & Standard model fails where clock succeeds; clock not unique \\
Fe/Mg vs $N$ & \textbf{[C]} & Textbook chemical evolution \\
Ionization plane & \textbf{[C]} & UV-background hardening \\
Velocity skewness & \textbf{[C]} & 7 bins; fragile \\
Straddling-pair coherence & \textbf{[B]} & z-immune by construction; single pass, unreplicated \\
Narrow-window doublet coherence & \textbf{[B]} & z-immune by construction; single pass, unreplicated \\
Entropy gradients & \textbf{[C]} & Structural; needs null modeling and replication \\
CaII low-z discrimination & \textbf{[A] inconclusive} & Pre-registered (exp\_18); no z-trend in CaII EW to discriminate on \\
Crossover at one Hartree (27.2 eV) & \textbf{[A] registered, untested} & Derived from M-R energy-scale machinery; registered \texttt{d9c77a81}; first measurement (exp\_19) blocked by cross-survey incomparability \\
CaII coupling = 0 & \textbf{[B]} & Measured with error bar (exp\_19): CI95 contains zero, unique among seven ions \\
\bottomrule
\end{longtable}

\begin{center}\rule{0.5\linewidth}{0.4pt}\end{center}

\section*{9. Conclusion}

The cascade clock makes first observational contact with quasar absorption spectroscopy. The contact established three things:

\textbf{A prediction that was falsified.} The oscillation was pre-registered, DFT-specific, and falsified. The failure scopes the clock's reach: it works at cosmological scales (S8, H0, JWST) but its non-monotonic features do not imprint on absorption line statistics at current precision. This is a statement about the instrument's leverage, not about PAC's validity.

\textbf{A partition worth keeping.} Photons carry conserved information (PAC: line ratios, topologically exact) and historical information (SEC: line shapes, epoch-dependent). This partition is DFT-specific — it predicts \emph{which} observables are in which class and \emph{why} (integer topology vs thermodynamic evolution). It survives regardless of the cascade clock's parameterization and could improve multi-line absorption analysis.

\textbf{A forward program, partly executed.} The cascade clock and its smooth mimics agree in the middle ($z = 1.5$–$4.5$) and diverge at the edges. The low-z edge has now been tested: the pre-registered CaII discrimination (exp\_18, commit \texttt{fbad01d1}) returned inconclusive — not because the models tie on a real trend, but because CaII equivalent widths carry no usable redshift trend at this sample size. The discriminating burden now rests on JWST at $z > 5$ and on low-z \emph{kinematic} observables. A derived ionization crossover energy from $\varphi$-scaling would provide a second non-degenerate test. These are the next experiments.

\begin{center}\rule{0.5\linewidth}{0.4pt}\end{center}

\section*{Data availability}

All code, data references, and pre-registration are at \texttt{https://github.com/dawnfield-institute/dawn-field-theory}, directory \texttt{foundational/experiments/midnight/}. Commit \texttt{193d1c8e} contains the pre-registered oscillation predictions; commit \texttt{fbad01d1} contains the pre-registered CaII low-z discrimination test (exp\_18).

\section*{References}

\begin{itemize}\tightlist
\item Anand, A. et al. (2021). SDSS DR16 MgII Absorber Catalog. MNRAS.
\item Monadi, R. et al. (2023). CIV Absorption Lines in SDSS DR12. Zenodo.
\item Davies, R. et al. (2023). XQR-30 Metal Absorber Catalog. MNRAS 521, 289.
\item Sardane, G.M., Rao, S.M., Turnshek, D.A. (2014). CaII absorbers in SDSS QSOs. MNRAS 444, 1747.
\item Groom, P. (2026). PACSeries (companion papers), Dawn Field Institute. Zenodo series concept DOI: 10.5281/zenodo.15783623 (resolves to the latest published version).
\item Webb, J.K. et al. (2011). Spatial variation of the fine structure constant. PRL 107, 191101.
\end{itemize}

\end{document}
